% Generated draft: requires mathematical and source review.
\documentclass[11pt]{article}
\usepackage[margin=25mm]{geometry}
\usepackage{fontspec}
\setmainfont{Latin Modern Roman}
\usepackage{amsmath,amssymb,mathtools,unicode-math}
\setmathfont{Latin Modern Math}
\usepackage{enumitem,needspace,xurl,hyperref}
\setlength{\parindent}{0pt}
\setlength{\parskip}{5pt}
\setlength{\emergencystretch}{3em}
\newcommand{\N}{\mathbb{N}}
\newcommand{\Z}{\mathbb{Z}}
\newcommand{\Q}{\mathbb{Q}}
\newcommand{\R}{\mathbb{R}}
\newcommand{\C}{\mathbb{C}}
\newcommand{\F}{\mathbb{F}}
\newcommand{\A}{\mathbb{A}}
\newcommand{\PP}{\mathbb{P}}
\newcommand{\E}{\mathbb{E}}
\newcommand{\eps}{\varepsilon}
\newcommand{\dd}{\,\mathrm{d}}
\begin{document}
\section*{UnsolvedMath ID: 30004252}
\textbf{Source problem code:} catalogue entry
\subsection{Definitions and mathematical statement}
Let \(G\) be a connected reductive \(p\)-adic group, let \(\mathcal B(G)\) denote its Bruhat--Tits building, and let \(x\in\mathcal B(G)\). The source also refers to a finite-length representation \(\tau(y,g)\) arising in the Mackey decomposition associated with a type, but does not specify the type, the inducing data, or the parameters \(y,g\).

The proposed geometric statement uses additional source-dependent objects \(H,G^0,\operatorname{proj}_0,\Omega_{t+},\Xi_t\), and \(s_0\), whose definitions are absent from the supplied record. If \(H\) acts on the building, its fixed-point set is
\[
\mathcal B(G)^H=\{u\in\mathcal B(G):hu=u\text{ for every }h\in H\}.
\]
For the intended auxiliary data, write
\[
A=\operatorname{proj}_0^{-1}\bigl(\mathcal B(G^0)\setminus\{x\}\bigr),\qquad
C_t=\Omega_{t+}\setminus\bigl(\Omega_{t+}\cap\Xi_t\bigr)
\quad(0<t<s_0).
\]
These expressions retain the source notation; they do not supply the missing definitions or identifications of buildings.

The proposed geometric assertion is that there exists a bounded convex neighborhood \(\Gamma\subseteq\mathcal B(G)\) of \(x\) such that
\[
\mathcal B(G)^H\subseteq\Gamma
\]
and, for every \(z\in\mathcal B(G)\setminus\Gamma\) and every geodesic segment \(\gamma\) from \(x\) to \(z\),
\[
\gamma\cap A\ne\varnothing
\quad\text{or}\quad
\exists\,t\in\mathbb R\;\bigl(0<t<s_0\ \text{ and }\ \gamma\cap C_t\ne\varnothing\bigr).
\]
Here the geodesic segment is the source's \([x,z]\); the existential parameter \(t\) may depend on the segment. The intended building, metric, topology, and admissible auxiliary data must be specified before this becomes a fully determined conjecture. In particular, containment in a bounded \(\Gamma\) requires \(\mathcal B(G)^H\) to be bounded, a property not justified by the supplied hypotheses.

The source further proposes that this geometric condition excludes additional types in \(\tau(y,g)\) far from \(x\). The relation between \(z\), the Mackey parameters \((y,g)\), and occurrence of a type is not supplied, and neither is a definition of the claimed small neighborhood. Consequently this representation-theoretic consequence cannot be formulated precisely from the record alone.
\subsection{Short English statement}
The proposed conjecture asks for a bounded convex neighborhood of \(x\), containing the \(H\)-fixed points, such that every geodesic from \(x\) to a point outside it intersects one of the specified obstruction sets. This is intended to localize the occurrence of types in Mackey components. Essential auxiliary definitions and the precise representation-theoretic consequence are missing, so this is a statement-only draft requiring mathematical review.
\subsection{Sources}
\begin{itemize}
\item Local UnsolvedMath problem registry, permanent ID \(30004252\), exact source problem code \texttt{catalogue entry}, title \emph{Localization of Types in Bruhat--Tits Buildings}; collection: Oberwolfach Reports Open Problems. The record supplies proposed year 2019 and names no proposer.
\item Supplied report citation: \emph{New Developments in Representation Theory of p-adic Groups} (2020), Oberwolfach Reports, pp.~2739--2819, DOI \texttt{10.4171/owr/2019/44}. \url{https://doi.org/10.4171/owr/2019/44}. The supplied curation identifies the assertion as Conjecture~2; the report text and its surrounding definitions were not supplied or independently consulted.
\item Supplied literature reference: \emph{Unicity of types for toral supercuspidals}. \url{https://arxiv.org/abs/1801.06721}. The record does not supply authors or a precise theorem supporting the assessment.
\item Literature assessment checked on \texttt{2026-08-22}: supplied dated curation, not independently verified. The record labels the problem open and states that the general anisotropic-torus case is poorly understood, with later results covering special unicity cases without establishing the full conjecture. Its background explicitly requests independent verification of time-sensitive status and cited claims. The record also labels the statement \texttt{corrected\_verified}, explaining that notation, typography, and sentence structure were repaired while preserving mathematical content; this label is not an independent mathematical verification of this draft.
\end{itemize}
\end{document}
